\documentclass[12pt]{article}
\usepackage[utf8]{inputenc}
\usepackage[T1]{fontenc}
\usepackage{hyperref}
\title{Muduga Language Resources for K-12 Curriculum Development}
\author{Compiled by Claude Sonnet 5 (claude-sonnet-5) for Andr\'as Kornai}
\date{}
\begin{document}
\maketitle
\section{Introduction}
Muduga (Mudugar, ISO 639-3 udg) is a Southern Dravidian language of the Attappady area of Palakkad district, Kerala, and nearby parts of Tamil Nadu. Wikipedia reports about 3,400 speakers (1991 census, so the figure is old). It is described as influenced by Kannada and Tulu and as mutually intelligible with Attapady Kurumba. Only a few resources turned up in a single search round.
\section{Primary Educational Resources}
No primers, textbooks or school programmes for Muduga were found.
\section{Linguistic and Reference Materials}
\begin{itemize}
\item Wikipedia article on Muduga, which lists as further reading a 1978 thesis by Rajendran, a 2017 descriptive grammar of Muduga and Kurumba by Shyam, Ravindran and Syama, and a 2022 article on centralized vowels by Arsenault and Abraham: \url{https://en.wikipedia.org/wiki/Muduga_language}
\item Related Wikipedia article on Attapady Kurumba: \url{https://en.wikipedia.org/wiki/Attapady_Kurumba_language}
\item Glottolog language page (entry reached through the ISO code pkr): \url{https://glottolog.org/resource/languoid/iso/pkr}
\end{itemize}
\section{Religious and Community Resources}
\begin{itemize}
\item Joshua Project profile of the Mudugar. Its language table conflicts with its own text, so treat it with caution: \url{https://joshuaproject.net/people_groups/17697}
\end{itemize}
\section{Digital and Audio Resources}
No audio, video or corpus resources were found.
\section{Curriculum Development Context}
The language appears to have no written educational materials and no known school use. A curriculum would have to start from documentation: an orthography decision, collection of texts and audio with the community, and the existing grammatical descriptions. Whether the sources above fit the udg code and the Muduga community should be checked first.
\end{document}